\documentclass[12pt,a4paper]{article}
\usepackage[T2A]{fontenc}
\usepackage[utf8]{inputenc}
\usepackage[russian]{babel}
\usepackage{amsmath,amssymb}
\usepackage[left=2.5cm,right=2cm,top=2cm,bottom=2cm]{geometry}
\usepackage{array}
% --- данные студента (поправьте при необходимости) ---
\newcommand{\Student}{Ярослав \rule{3cm}{0.4pt}}
\newcommand{\Group}{гр. \rule{1.5cm}{0.4pt}}

\begin{document}
\begin{titlepage}
\centering
Московский государственный университет имени\\ М.\,В.~Ломоносова\\[8pt]
Механико-математический факультет
\vfill
{\Large Отчёт по задаче №1 практикума на ЭВМ}
\vfill
\begin{flushright}\Student\\ \Group\end{flushright}
\vfill
Москва\\ 2026
\end{titlepage}

\section{Постановка задачи (№1.17)}
Для каждого действительного значения параметра $\alpha$ требуется численно найти решение уравнения
\begin{equation}\label{eq:main}
1+e^{-2x^2}=\alpha^2e^{-2x}.
\end{equation}
Решение реализуем методом итераций с начальным приближением $x_0=0$. Требуемая точность
(по модулю) равна $\varepsilon=10^{-12}$.

\section{Разрешимость уравнения}
Левая часть \eqref{eq:main} не меньше $1$, поэтому при $\alpha=0$ решений нет.
Пусть $\alpha\ne0$. Обе части \eqref{eq:main} положительны; логарифмируя, получаем равносильное уравнение
\begin{equation}\label{eq:log}
F(x):=2x+\ln\bigl(1+e^{-2x^2}\bigr)=2\ln|\alpha| .
\end{equation}
Уравнение зависит от $\alpha$ только через $|\alpha|$, поэтому решения при $\alpha$ и $-\alpha$ совпадают.

\section{Приведение к виду неподвижной точки}
Выразим из \eqref{eq:log} слагаемое $2x$. Получим равносильное уравнение
\begin{equation}\label{eq:fixed}
x=G_\alpha(x),\qquad G_\alpha(x)=\ln|\alpha|-\frac12\ln\bigl(1+e^{-2x^2}\bigr).
\end{equation}
Следовательно, решение исходного уравнения является неподвижной точкой отображения $G_\alpha\colon\mathbb R\to\mathbb R$.

\section{Доказательство сжимаемости отображения}
Функция $G_\alpha$ дифференцируема на всей прямой, и
\[
G_\alpha'(x)=-\frac12\cdot\frac{-4x\,e^{-2x^2}}{1+e^{-2x^2}}=\frac{2x\,e^{-2x^2}}{1+e^{-2x^2}}=\frac{2x}{e^{2x^2}+1}.
\]
Параметр $\alpha$ в производную не входит. Используя неравенство $e^{t}\ge1+t$ при $t=2x^2$, а затем $2|x|\le1+x^2$, получаем
\begin{equation}\label{eq:k}
|G_\alpha'(x)|=\frac{2|x|}{e^{2x^2}+1}\le\frac{2|x|}{2+2x^2}=\frac{|x|}{1+x^2}\le\frac12=k<1
\qquad\forall x\in\mathbb R.
\end{equation}
По формуле Лагранжа $|G_\alpha(u)-G_\alpha(v)|=|G_\alpha'(\xi)|\,|u-v|\le k|u-v|$, т.\,е.\
$G_\alpha$ является сжимающим отображением полного пространства $\mathbb R$. По теореме Банаха
при каждом $\alpha\ne0$ существует единственная неподвижная точка $x^*$, т.\,е.\ уравнение
\eqref{eq:main} имеет ровно одно решение, и итерации $x_{n+1}=G_\alpha(x_n)$ сходятся к нему
при любом начальном приближении.

\section{Итерационный алгоритм и оценка погрешности}
Положим $x_0=0$, $x_{n+1}=G_\alpha(x_n)$. Из неравенства треугольника и свойства сжатия
\[
|x_n-x^*|\le|x_n-x_{n+1}|+|G_\alpha(x_n)-G_\alpha(x^*)|\le|x_n-x_{n+1}|+k|x_n-x^*|,
\]
откуда
\[
|x_n-x^*|\le\frac{|x_{n+1}-x_n|}{1-k}.
\]
При $k=1/2$ условие достижения точности $\varepsilon=10^{-12}$ имеет вид
\begin{equation}\label{eq:stop}
\frac{|x_{n+1}-x_n|}{1-k}\le\varepsilon\iff|x_{n+1}-x_n|\le\frac{10^{-12}}{2}=5\cdot10^{-13}.
\end{equation}

\section{Результаты работы программы}
Программа находит решение уравнения для заданного вещественного значения параметра $\alpha$,
если оно передано единственным аргументом командной строки. Если переданы два целых числа
$a<b$, программа находит решения для всех целых $\alpha$ из отрезка $[a,b]$, включая его концы.
При $\alpha=0$ программа сообщает, что решений нет.

В таблице приведены результаты работы программы для целых $\alpha$ от $-10$ до $10$
($n$ — порядковый номер приближения, удовлетворяющего искомой точности).

\begin{center}
\renewcommand{\arraystretch}{1.12}
\begin{tabular}{r r r}
\hline
$\alpha$ & \multicolumn{1}{c}{$x_n$} & $n$\\
\hline
$-10$ & $2.302572680928053$ & $5$ \\
$-9$ & $2.197192536178525$ & $5$ \\
$-8$ & $2.079353760117693$ & $5$ \\
$-7$ & $1.945652659970291$ & $6$ \\
$-6$ & $1.790941682762951$ & $7$ \\
$-5$ & $1.606581570647285$ & $9$ \\
$-4$ & $1.375026870307486$ & $12$ \\
$-3$ & $1.045341592097919$ & $19$ \\
$-2$ & $0.430809418616581$ & $26$ \\
$-1$ & $-0.302822270174649$ & $21$ \\
$0$ & \multicolumn{1}{c}{решений нет} & --- \\
$1$ & $-0.302822270174649$ & $21$ \\
$2$ & $0.430809418616581$ & $26$ \\
$3$ & $1.045341592097919$ & $19$ \\
$4$ & $1.375026870307486$ & $12$ \\
$5$ & $1.606581570647285$ & $9$ \\
$6$ & $1.790941682762951$ & $7$ \\
$7$ & $1.945652659970291$ & $6$ \\
$8$ & $2.079353760117693$ & $5$ \\
$9$ & $2.197192536178525$ & $5$ \\
$10$ & $2.302572680928053$ & $5$ \\
\hline
\end{tabular}
\end{center}
\end{document}
